\documentclass{article}

%A mettre tout le temps
\usepackage{lcpackage}

\begin{document}
\pagestyle{fancy}
\fancyhf{}
\fancyhead[L]{{\fontfamily{phv}\selectfont Matière}}
\fancyhead[C]{{\fontfamily{phv}\selectfont $\cdots$ Titre - Nom $\cdots$}}
\fancyhead[R]{{\fontfamily{phv}\selectfont Page \thepage /\pageref{LastPage}}}
{\fontfamily{phv}\selectfont
\begin{center}
\Huge{Titre} %Titre à changer
\end{center}
Voici un document de test:\\
\begin{definition}{Titre definition}
Voici du texte dans lequel je vais inscrire une formule:
\formulev{$Formule=cool$}
\end{definition}
\noindent Si je veux faire un programme Python:
\begin{prog}{Nom du prog que je ne suis pas obligé de mettre}
def fonction(n):
	a=0
	for k in range(n):
	   a+=1
	return(a)
\end{prog}
\noindent Pour faire du SQL:

\begin{progSQL}{}
SELECT prop1
FROM prop
WHERE prop1 >=2;
\end{progSQL}
\noindent Lorsque je dois faire des démonstrations ou exercices de maths:\\
\xBox{Soit $x\in\R$}{
Début de la démo\\
\Donc Je peux faire apparaitre de nouvelles hypothèses\\
\xBox{$\boxed{\Rightarrow}$ Supposons un truc}{
\Mq{Montrons que machin}\\
\But{Notre but sera de faire ça}
}{et ce en supposant le truc}
}{et ce en supposant $x \in \R$}\\
Je peux très bien faire des \pg{$\backslash$begin$\{$multicols$\}$}:
\begin{multicols}{2}
\begin{méthode}{Comment rédiger une méthode}
Une méthode s'écrit toujours avec des étapes complètes. Il faut tout détailler, y mettre les formules. Des fois cela nécessite des sous étapes. Ce n'est pas la peine d'écrire de longue phrases. On peut très bien écrire une phrase commençant par \textbf{un verbe à l'infinitif ou des On fait ça} puis mettre une action.

On peut aussi mettre des \hyperref[methode1]{\textcolor{myorange}{\textbf{liens}}} vers des définitions ou autres trucs:
\end{méthode}\columnbreak
\begin{méthode}[label=methode1]{Exemple d'une méthode}\vspace*{-7pt}
\begin{itemize}
\item Pour un $M$ qlq, on pose le syst donné et on trouve une équation du plan
\item On montre que $M_0 \in \mathscr C$
\item Méthode gradient:
\iitem Calcul du $\nabla f(x,y)$
\iitem On remplace par les cos de $M_0$
\iitem On a le vecteur normal d'où la tangente
\item Méthode DL:
\iitem On pose le paramètre $t=3$ ($t$ qui permet de retrouver le point via le système)
\iitem On pose $t=3+h$ + changement dans le système
\iitem DL associé à chaque équation
\iiitem Pour la première on simplifie et on rajoute un $o(h)$
\iiitem Pour la seconde, Id remarquable et on enlève le terme en $h^2$ et on rajoute un $o(h)$
\iitem On obtient alors les coordonnées d'un vecteur directeur pour notre tangente !
\end{itemize}
\end{méthode}
\end{multicols}
\begin{demomethode}{Titre de la démo}{
\begin{itemize}
\item Méthode pas à pas de la démo.
\item Puis
\end{itemize}
}
Contenu démo (possibilité de couper des pages)

\begin{multicols}{2}
    On peut aussi mettre des colonnes internes.\columnbreak\\
    Colonne 2 :)
\end{multicols}

\begin{progSQL}{Implémenter des démo-méthodes}
\begin{demomethode}{Titre de la démo}{
Contenu méthode
}
Contenu démo
\end{demomethode}
\end{progSQL}
\end{demomethode}
\begin{exemplemethode}{Titre de l'exemple}{
}
Il existe pareil mais avec des \pg{exemplemethode}.
\end{exemplemethode}
\end{document}